% ==============================================================================
% T0 Theory / FFGFT: Document 284 (Chapter Content, EN)
% Title: A formula-level reading of the Spiral-Time neurodynamic model
% Author: Johann Pascher
% Date: 19 June 2026
% Purpose: Reproducible, formula-level test of Krueger/Feeney/Wende
%          (Medinformatics 2026); the foundational time-axis issue; the FFGFT
%          compactification / conversion-as-measurement reading.
% References: Doc. 190, 210, 249, 263, 270, 282, 283.
% Scripts: test_algebra.py, test_deltaphi.py, test_dynamics.py,
%          test_compactification.py, make_figures.py
% ==============================================================================

\section*{Abstract}
This document is a reproducible, formula-level reading of \emph{Information-Theoretic
Modeling of Neural Coherence via Triadic Spiral-Time Dynamics: A Framework for
Neurodynamic Collapse} (Kr\"uger, Feeney \& Wende, \emph{Medinformatics} 2026,
DOI 10.47852/bonviewMEDIN62029043). Every verdict is produced by a short runnable
script. The reading separates what holds (an interpretable phase-deviation feature;
genuine octonion non-associativity) from what is mislabeled (the triadic algebra is
quaternionic, not bicomplex), internally inconsistent (two incompatible definitions of
$\Delta\Phi$), not derived (the $0.40$ threshold; the claimed spirals and limit cycles),
or not reproducible from the paper (the reported $r\approx0.96$). Underneath these sits
\emph{one} foundational issue: the operator $\psi(t)=t+i\varphi(t)+j\chi(t)$ is built on
physical, linear time as its base axis -- the decompactified, post-measurement coordinate
-- so its memory is generic by construction. The FFGFT reading makes this precise: the
system is compact and carries memory. By Doc.~270 the decompactification to SI time is a
mode-preserving embedding (lossless), and FFGFT does carry framework-specific SI
observables (masses/Koide, CHSH$\sim\xi$, $D_f$, $L_0$; Doc.~270/249); what removes the
recurrence is not the SI conversion but the finite, non-periodic observation window plus
decoherence. The sharp, narrow consequence: the coherence-sector recurrence is not
\emph{resolvable} on a warm, finite-window SI-time signal such as EEG. This is a reviewer note addressed to the authors; the
intent is the falsification-forward standard the HLV/HICE side already applies.

% ==============================================================================
\section{Occasion and method}
% ==============================================================================
The Medinformatics paper applies the HLV Spiral-Time operator to neural coherence and
proposes the Information--Coherence--Energy (ICE) stability principle and a
$\Delta\Phi$-based ``neurodynamic collapse'' boundary. Rather than argue from
impression, each mathematical claim here is checked by computation. The scripts are
\texttt{test\_algebra.py}, \texttt{test\_deltaphi.py}, \texttt{test\_dynamics.py},
\texttt{test\_compactification.py} and \texttt{make\_figures.py}; fixed seeds throughout.
Findings split into: \emph{holds}, \emph{mislabeled}, \emph{internally inconsistent},
\emph{not derived}, \emph{not reproducible from the paper}, and \emph{speculative
premise}. One issue is foundational -- the time-axis the operator is built on -- and is
taken first, because findings on $\Delta\Phi$ and on the dynamics are its symptoms.

\begin{center}\small
\renewcommand{\arraystretch}{1.25}
\begin{tabular}{@{}c p{4.1cm} p{1.7cm} p{5.0cm}@{}}
\toprule
\# & Claim & Section & Status\\
\midrule
0 & Base time-axis $\psi(t)=t+i\varphi(t)+j\chi(t)$ on physical linear time & 2.1, 2.2.1 & Foundational -- operator on the post-measurement axis (subsumes 3 \& 5)\\
1 & Triadic time ``reduced bicomplex/octonionic plane'' & 2.1 & Mislabeled -- it is the quaternions\\
2 & Octonionic operator; non-assoc.\ ``explains psychology'' & 2.10 & Math holds, claim overreaches\\
3 & $\Delta\Phi$ instability functional & 2.6 vs 12/15 & Internally inconsistent (two definitions)\\
4 & Critical $\Delta\Phi_c\approx0.40$, ``not a free parameter'' & 2.5, 2.6.2 & Not derived (free parameter)\\
5 & Healthy spirals / trauma limit cycles & 2.2.3, 2.9 & Not derived from the equations\\
6 & $r\approx0.96$, $\Delta\Phi_c\approx0.40$ ``across subjects'' & 3.4 & Not reproducible from the paper\\
7 & Axonal biphoton entanglement coupling $\Gamma_{\text{axon}}$ & 2.7--2.8 & Speculative premise as established\\
\bottomrule
\end{tabular}
\end{center}

% ==============================================================================
\section{The base time-axis (foundational)}
% ==============================================================================
This is the most consequential choice in the paper, and the one under which findings 3
and 5 are downstream symptoms. The Spiral-Time operator
\begin{equation}
\psi(t) = t + i\,\varphi(t) + j\,\chi(t)
\end{equation}
takes physical, linear time $t$ as its scalar base and writes $\varphi,\chi$ as
functions of that same $t$. The paper is explicit that $t$ is physical: $U_1$ is defined
as ``linear physical time,'' tracking external chronology (\S2.2.1). That fixes the
representation in which every later memory and holonomy claim is made.

\paragraph{Unit-independent.} For $t + i\varphi + j\chi$ to be one algebraic object the
three components must be the same kind of quantity. The paper asserts dimensionlessness
(\S2.1.2) yet simultaneously calls $t$ physical time and chronology (\S2.2.1); $\varphi$
and $\chi$ are internal phase / identity coordinates. Adding a physical time to internal
phases in one quaternionic sum is dimensionally incoherent unless all three are first
reduced to a common dimensionless measure. (Since $i,j$ anticommute the object is
quaternionic -- $t$ is the real part of a quaternion whose cross-term $ij=k$ is never
tracked; see \S\ref{sec:alg}.)

\paragraph{Deeper.} A framework-specific memory lives on the compact internal time-circle
$S^1$: a discrete spectrum, an exact recurrence (in FFGFT the coherence-sector back-flow
$5.125$, Doc.~282). Writing $\varphi(t),\chi(t)$ as single-valued functions of physical
linear $t$ flattens $S^1$ onto $\mathbb{R}$ -- it performs the projection
$S^1\!\to\!\mathbb{R}$ before any dynamics, and a single-valued function over linear time
has no compact periodicity, so the recurrence is not representable on this axis at all.
By the conversion-as-measurement principle (\S\ref{sec:ffgft}), committing to physical
time is itself the disturbing, information-discarding projection -- no apparatus needed.
The notation has silently taken that step, hiding the one operation -- the
conversion/measurement from compact internal time to physical SI time -- that decides
whether a memory claim is framework-specific or generic.

\paragraph{Consequence.} The operator is, by construction, post-measurement: whatever
memory it carries on physical (EEG) time is generic. This is reproducible
(\texttt{test\_compactification.py}): the discrete-spectrum return signal has back-flow
$11.5$ on the compact horizon, which a finite physical read-out window drives to $0.000$,
while the six lines never resolve ($2\pi/T_{\text{win}}\gg$ line spacing $0.091$) -- the
read-out is a non-invertible projection. The same root explains the two downstream
findings: a Berry-phase holonomy (Eq.~12/15, \S\ref{sec:dphi}) needs a closed loop -- a
compact cycle -- and a limit cycle (\S\ref{sec:dyn}) needs non-gradient closed orbits;
neither is supported on a linear physical-time axis.

\begin{figure}[h]\centering
\includegraphics[width=0.95\linewidth]{fig_compactification.png}
\caption{Read-out as a non-invertible projection. Left: on the compact representation the
discrete spectrum produces revivals (back-flow $\gg0$); the red band marks a finite
physical read-out window. Right: inside the window the observable decays without revival
and the discrete lines are unresolved.}
\end{figure}

\paragraph{Constructive fix.} Define the operator on the compact, dimensionless internal
time (cyclic), where the discrete spectrum and recurrence live, and let physical SI time
enter only through an explicit, declared conversion -- which is precisely what projects
out the framework-specific part. Either the operator stays internal/compact (no direct
EEG claim) or it is physical/SI (generic memory). $\psi(t)=t+i\varphi(t)+j\chi(t)$ tries
to be both. The minimal change is to declare two distinct times and write the conversion
between them explicitly as the measurement it is.

% ==============================================================================
\section{Triadic time algebra -- mislabeled}\label{sec:alg}
% ==============================================================================
The stated relations $i^2=j^2=-1,\ ij=-ji$ are exactly the imaginary quaternions: setting
$k:=ij$ gives $k^2=-1$, and $\{1,i,j,k\}$ closes into $\mathbb{H}$, which is associative
and non-commutative (\texttt{test\_algebra.py}, A). A \emph{bicomplex} plane requires the
two imaginary units to commute ($ij=ji$); here $ij=-ji$, so the algebra is \emph{not}
bicomplex. A three-term operator $t+i\varphi+j\chi$ lives inside a quaternion
sub-algebra; octonions are not required for it. The correct name is quaternionic.

% ==============================================================================
\section{Octonionic structure -- math holds, interpretation overreaches}
% ==============================================================================
Octonion non-associativity is real: the associator $(e_1e_2)e_4-e_1(e_2e_4)=2e_7\neq0$
(Cayley--Dickson), while associators on a quaternionic triple vanish
(\texttt{test\_algebra.py}, B). That is correct mathematics. The further step -- that this
``directly matches empirical psychology'' and is ``the first physics-based explanation for
psychological nonlinearity'' -- is an interpretive analogy, not a derivation: no
psychological quantity is computed from or tested against the algebra, and the eight basis
directions are assigned meanings by stipulation.

% ==============================================================================
\section{$\Delta\Phi$ -- two incompatible definitions under one symbol}\label{sec:dphi}
% ==============================================================================
\begin{itemize}
\item \textbf{Definition 1 (\S2.6):} $\Delta\Phi=\alpha|\Delta S|+\beta|\Delta I|+\gamma|\Delta C|$,
convex weights, each component in $[0,1]$ -- a unitless weighted-$L^1$ index in $[0,1]$.
It is linear in each component, so it has no curvature and no intrinsic bifurcation; the
level set $\Delta\Phi=c$ is a flat polytope facet for every $c$.
\item \textbf{Definition 2 (Eq.~12, 15):} $\Delta\Phi=\oint(\nabla\times A)\cdot da
\approx \mathrm{Im}\oint\langle\psi|\nabla\theta|\psi\rangle\,d\theta$ -- a Berry-phase
holonomy, an angle in radians ranging over $0\ldots2\pi$, with no reason to lie in $[0,1]$
or near $0.40$ (a loop with $|\gamma|=0.40$\,rad is a tiny $\sim$\,29$^\circ$ cone).
\end{itemize}
A normalized deviation index and a radian holonomy are different mathematical objects.
Calling both ``$\Delta\Phi$'' and assigning both the same critical value $0.40$ is an
equivocation (\texttt{test\_deltaphi.py}, C/D).

\begin{figure}[h]\centering
\includegraphics[width=0.72\linewidth]{fig_deltaphi.png}
\caption{The two $\Delta\Phi$ definitions are incompatible objects: a radian holonomy
($0\ldots2\pi$) versus a unitless index ($0\ldots1$).}
\end{figure}

% ==============================================================================
\section{The $0.40$ threshold -- not derived}
% ==============================================================================
\S2.5 states the value is ``not imposed as a free tuning parameter'' and that the operator
structure ``enforces'' a bounded instability surface. But the $L^1$ functional has no
intrinsic critical point (it is linear); the only structural statement available is the
trivial bound $0\le\Delta\Phi\le1$ and the symmetric reference $\Delta\Phi=d$ (so $1/3$ at
$d=1/3$). \S2.6.2 itself attributes the $0.40$-vs-$0.333$ gap to ``dataset-specific
weighting asymmetry $\ldots$ and normalization,'' and at equal deviation $d=0.40$ every
convex weight set returns $0.40$ (\texttt{test\_deltaphi.py}, C/E). So $0.40$ is an
empirical/contour value, not a structurally enforced bifurcation.

% ==============================================================================
\section{Spirals and trauma limit cycles -- not derived}\label{sec:dyn}
% ==============================================================================
The only dynamical mechanism given with explicit potentials (\S2.4) is the gradient-flow
part of the master equation, $\partial_t\Psi=-\nabla V_{\text{DLHR}}(\Psi)$. For any
gradient flow,
\begin{equation}
\frac{dV}{dt}=\nabla V\cdot\frac{d\Psi}{dt}=-|\nabla V|^2\le0,
\end{equation}
so $V$ is a strict Lyapunov function and the system cannot have periodic orbits, limit
cycles, or sustained spirals -- it relaxes to fixed points (theorem). Numerically
(\texttt{test\_dynamics.py}, F): across 40 trajectories $V(t)$ is monotonically decreasing
(0 violations) and 0 closed orbits are found. The phenomenology ``healthy cognition
$\to$ stable inward spirals,'' ``trauma loops $\to$ limit cycles'' therefore does not
follow from the specified dynamics; it would require the advection terms
$\lambda_\varphi\partial_\varphi,\ \lambda_\chi\partial_\chi$ or $W_{\text{oct}}$, which
are never given concretely. (The cubic $M^3$ term also makes $V$ unbounded below.)

\begin{figure}[h]\centering
\includegraphics[width=0.95\linewidth]{fig_dynamics.png}
\caption{Gradient flow with the paper's DLHR potentials: trajectories relax to fixed
points (no closed orbits, no sustained spiral) and $V(t)$ decreases monotonically, so
limit cycles are impossible.}
\end{figure}

% ==============================================================================
\section{Empirical claims -- not reproducible from the paper}
% ==============================================================================
The reported $r\approx0.96$ and $\Delta\Phi_c\approx0.40$ ``across subjects'' (\S3.4)
cannot be recomputed from the paper, which gives no data table. Reproduction needs the
external dataset (Zenodo 10.5281/zenodo.17990696) and the validation notebook. A fair
test would require the dataset plus a pre-registered protocol with the nulls the authors
already use on the HICE side (time-shuffle, present-state-matched, phase-randomized).

% ==============================================================================
\section{Axonal biphoton entanglement -- speculative premise stated as established}
% ==============================================================================
$\Gamma_{\text{axon}}=g\,e^{i\varphi}E_{12}+g'e^{j\chi}E_{23}$ (\S2.8) is built on
``experimental findings'' of biphoton entanglement and long-range coherence in myelinated
axons, presented as settled research. Entanglement in warm, wet, myelinated axons is not
established physics. This is not a formula error, but it is a load-bearing empirical
premise asserted as fact and should be flagged as hypothesis.

% ==============================================================================
\section{A defensible core}
% ==============================================================================
Stripped of the superstructure there is a sound, modest object:
$\Delta\Phi=|\omega(t)-\omega_{\text{memory}}(t)|$ -- the deviation of instantaneous phase
velocity from a short-term phase-memory baseline -- is an interpretable, lightweight
time-series instability feature, and as a signal-analysis operator (the framing of the
companion respiratory paper) it is legitimate. The problems are confined to the
HLV/octonion/consciousness superstructure, the $\Delta\Phi$ redefinition, and the
un-derived threshold and dynamics.

% ==============================================================================
\section{An FFGFT reading: conversion as measurement}\label{sec:ffgft}
% ==============================================================================
Temporal recursion does carry memory, and it maps onto FFGFT's own structure (the
recursion operator $\mathcal{R}$; the memory kernel as the Fourier transform of the
connection-Laplacian spectral density, Doc.~282). The following is offered as a
constructive bridge and an explicit conjecture, not a result.

FFGFT's memory has two faces. On the overdamped / relaxation observable the kernel is a
positive multi-exponential, and a generic biexponential reproduces it (residual
$\sim10^{-3}$): non-identifiable, Stufe-0. On the coherence observable
$c(t)=e^{-\Gamma(t)}$ -- its home sector -- the discrete spectrum produces a genuine
revival: back-flow $5.125$, which generic decaying kernels cannot reproduce (back-flow
$0$). So the memory is discriminating, but only on the coherence sector and only when
\emph{both} a discrete spectrum \emph{and} preserved coherence (small $\Gamma$) hold.

FFGFT is fundamentally compact ($T^4$, discrete spectrum), and in that dimensionless
representation the system always has memory. It is important to locate the loss precisely,
because the conversion to SI is \emph{not} where it happens. By Doc.~270 the Type~I duality
decompactification $S^1_m\!\to\!\mathbb{R}_t$ (governed by $\tilde{T}\cdot m=1$) is a
\emph{mode-preserving embedding} -- lossless; the periodicity merely stops being directly
visible and survives as the discrete spectrum and as physical observables (the revival time
$T_\text{rev}$, the CHSH deviation $\sim\xi$). A bare units rescaling (by $\hbar,c$) is
likewise invertible and loses nothing. The genuinely lossy operations are three and
distinct: the \emph{reverse} compactification $\mathbb{R}\!\to\!S^1$ (non-periodic parts do
not descend), the spatial projections $D3\!\to\!D0$ (Doc.~270), and -- for any read-out --
the restriction to a \emph{finite, non-periodic observation window plus decoherence}. It is
this last step, not the unit conversion, that keeps the local short-time behaviour and
discards the global recurrence, exactly as a measurement selects a basis and drops the
coherences (Test~G). A physical measurement is the same window-and-decoherence projection
at its strong-coupling end. This is consistent with what the corpus already records: the
vacuum energy is present compactified, visible decompactified, with $c,\hbar c$ as SI
conversion factors (Doc.~263) -- the decompactified view is not less true, only differently
resolved.

The reading therefore both supports the modest core (recursion does carry memory) and
bounds the claim precisely. The framework-specific signature does survive into SI in the
right channels -- mass reading, CHSH$\sim\xi$, $D_f$, $L_0$ (Doc.~270/249). What does
\emph{not} survive is the coherence-sector recurrence read on a warm, finite-window SI-time
observable: there the window and decoherence have already discarded the global recurrence,
so the very observable the clinical sections use -- SI-time EEG -- is post-measurement in
exactly that sense, and is the regime where the FFGFT harness (Doc.~282/283) shows the
memory to be generic, not a fingerprint. This is also why the base-time-axis choice (\S3)
is decisive: the operator is written on that finite, post-measurement axis from the start.

% ==============================================================================
\section{Reproducibility}
% ==============================================================================
\begin{itemize}
\item \texttt{test\_algebra.py} -- quaternion closure + bicomplex refutation (A); octonion associator (B).
\item \texttt{test\_deltaphi.py} -- $L^1$ bound/linearity (C); Berry-phase range (D); $0.40$ as free knob (E).
\item \texttt{test\_dynamics.py} -- gradient-flow Lyapunov / no-limit-cycle test (F).
\item \texttt{test\_compactification.py} -- read-out as a non-invertible projection: compact back-flow vs windowed (G).
\item \texttt{make\_figures.py} -- the figures used above.
\end{itemize}
The figures \texttt{fig\_deltaphi.png}, \texttt{fig\_dynamics.png} and
\texttt{fig\_compactification.png} are produced by these scripts; all use fixed seeds.
